\documentclass[10pt,a4paper]{amsart}
\usepackage[margin=19mm]{geometry}
\usepackage{amsmath,amssymb,mathtools}
\usepackage[scr=rsfs,cal=euler]{mathalfa}
\usepackage{xfrac}
\usepackage[unicode,hidelinks,bookmarks=false]{hyperref}
\newcommand{\ie}{i.e.\ }
\newcommand{\cf}{cf.\ }
\newcommand{\resp}{respectively\ }
\newcommand{\Subsection}{\subsection}
\providecommand{\nobreakdash}{\nobreak\hbox{-}}
\setlength{\emergencystretch}{2em}
\multlinegap0pt
\usepackage{bm}
\usepackage{bbm}
\usepackage{dsfont}
\usepackage{xspace}
\usepackage{cancel}
\usepackage{mathtools}
\usepackage{amsthm}
\makeatletter
\@ifundefined{thm}{\newtheorem{thm}{Theorem}}{}
\makeatother
\title[Adjacency-diametrical matrix of a graph]{Adjacency-diametrical matrix of a graph}
\author{S. P. Leka Amruthavarshini, R. Rajkumar}
\date{Source: arXiv:2601.01193v1 (CC BY 4.0)}
\begin{document}
\maketitle

\noindent\emph{Source and licence.} Adjacency-diametrical matrix of a graph. Version arXiv:2601.01193v1 (CC BY 4.0), distributed under the Creative Commons Attribution 4.0 International licence, as recorded in the arXiv licence metadata for this exact version and shown on the abstract page https://arxiv.org/abs/2601.01193v1. Full rights notice: licenses/arxiv-2601.01193-CC-BY-4.0.txt.\par\smallskip

\noindent\textit{Editorial scope.} Counted unit: the Thm whose source label is \texttt{None} in the article (source0.tex lines 243--246, source label \texttt{None}), delivered together with the article's own complete original proof of it (source0.tex lines 247--327, one \texttt{proof} environment of 81 lines). No mathematics is written, repaired or re-derived: statement and proof are the article's own text line for line. Required context retained uncounted: none. Every definition, display and label of those lines is retained unchanged. Omitted from this package and not redistributed: 1--242, 328--974. Nothing inside a retained excerpt is omitted or altered: every retained body is delivered exactly as the article's lines are written, so no omission receipt of any kind accompanies this package. Citations: the retained excerpts carry no citation command at all, so no bibliography entry of the article is transcribed and no reference is redistributed. Reference adaptation: none was needed and none was made; every reference inside the delivered unit resolves inside it, so no unresolved reference or citation is produced. Printed slips kept literal: no printed word, symbol, hypothesis or number of the counted statement or of its proof was altered. Layout and class: the article's own class is \texttt{article} with packages including amssymb, amsmath, amsthm, array, hyperref, caption, subcaption, float, longtable, multirow, accents, bm, url, geometry; the wrapper is the lane's pinned \texttt{amsart} editorial wrapper at 10pt on A4, which restates the theorem-like environments the retained text needs and adds the article's own macro definitions that the retained text needs (none), with extra wrapper packages (bm, bbm, dsfont, xspace, cancel, mathtools). The delivered numbering is the unit's own. Assets: no retained span contains a figure environment, a graphics-inclusion command, a PDF-page inclusion, a table or a TikZ picture; no image file is copied into this package, the package declares no asset dependency and no image file is redistributed with it. Licence: version arXiv:2601.01193v1 of this article is distributed under the Creative Commons Attribution 4.0 International licence, as recorded in the arXiv licence metadata for this exact version and shown on the abstract page https://arxiv.org/abs/2601.01193v1. A full rights notice is delivered at licenses/arxiv-2601.01193-CC-BY-4.0.txt. Characters: every retained excerpt is pure ASCII, so the delivered engine, XeLaTeX with its default Latin Modern text font, reports no missing character.
	\begin{thm}
	The characteristic polynomial of $AD(S_{n_1,n_2})$ is 
		\[x^{n_1+n_2-4}\left(x^4-\left(9n_1n_2-8n_1-8n_2+8\right)x^2+4\left(n_1n_2-n_1-n_2+1\right)\right).\]
	\end{thm}

	\begin{proof}
	By suitably arranging the vertices of $S_{n_1,n_2}$, it can be seen that
		\[AD(S_{n_1,n_2})=\begin{bmatrix}
		0 & 1 & J_{1 \times (n_1-1)} & \bm{0}_{1 \times (n_2-1)}\\
		1 & 0 & \bm{0}_{1 \times (n_1 - 1)} & J_{1 \times (n_2 -1)}\\

		J_{(n_1 - 1) \times 1} & \bm{0}_{(n_1-1) \times 1} & \bm{0}_{(n_1-1) \times (n_1-1)} & 3J_{(n_1-1) \times (n_2-1)}\\
		\bm{0}_{(n_2-1) \times 1} & J_{(n_2-1) \times 1} & 3J_{(n_2-1) \times (n_1-1)} & \bm{0}_{(n_2-1) \times (n_2-1)}
		\end{bmatrix}.\]
		

		
		Since each row of $\bm{0}_{(n_1-1) \times (n_1-1)}$ (resp. $\bm{0}_{(n_2-1) \times (n_2-1)}$) has equal row sum $0$, the vector $\bm{1}_{(n_1-1) \times 1}$ (resp. $\bm{1}_{(n_2-1) \times 1}$) is an eigenvector corresponding to the eigenvalue $0$. Assume that the remaining eigenvectors $X_i$ for $i=2,3,\dots,n_1-1$ (resp. $Y_j$ for $j=2,3,\dots,n_2-1$) are chosen orthogonal to $\bm{1}_{(n_1-1) \times 1}$ (resp. $\bm{1}_{(n_2-1) \times 1}$). 
		Then $ \left[
		0 \quad 0 \quad X^T_i \quad \bm{0}_{1 \times (n_2-1)}
		\right]^T$ 
		$\big(\text{resp.} \left[
			0 \quad  0 \quad \bm{0}_{1 \times (n_1-1)} \quad Y^T_j
		\right]^T\big)$
       is an eigenvector of $AD(S_{n_1,n_2})$ corresponding to the eigenvalue $0$.
			
		
		This implies that for $i=2,3,\dots,n_1-1,$ and for $j=2,3,\dots,n_2-1$, the vectors
		$\left[
		0 \quad 0 \quad X_i \quad \bm{0}_{1 \times (n_2-1)}
		\right]^T$ and 
		$\left[
	 0 \quad 0 \quad \bm{0}_{1 \times (n_1-1)} \quad Y_j
	 \right]^T$ are eigenvectors of $AD(S_{n_1,n_2})$ corresponding to the eigenvalues $0$ with multiplicity $n_1-2$ and $0$ with multiplicity $n_2-2$ respectively. Therefore $0$ is an eigenvalue of $AD(S_{n_1,n_2})$ with multiplicity $n_1+n_2-4$. In this way we obtained $n_1+n_2-4$ eigenvalues. 		  
	 All these eigenvectors are orthogonal to $\left[
	 1 \quad 0 \quad \bm{0}_{1 \times (n_1-1)} \quad \bm{0}_{1 \times (n_2-1)}
	 \right]^T$,
	 $\left[
	 0 \quad 1 \quad \bm{0}_{1 \times (n_1-1)} \quad \bm{0}_{1 \times (n_2-1)}
	\right]^T$,
	  $\left[
	 0 \quad 0 \quad J_{1 \times (n_1-1)} \quad \bm{0}_{1 \times (n_2-1)}
	\right]^T$ 	 	 and 	 	 $\left[
	 	0 \quad 0 \quad \bm{0}_{1 \times (n_1-1)} \quad J_{1 \times (n_2-1)}
	\right]^T$. 
	 
	
	 
	 Hence, the space spanned by these four vectors and the space spanned by the remaining four eigenvectors of $AD(S_{n_1,n_2})$ are same. Therefore, the remaining eigenvectors of $AD(S_{n_1,n_2})$ are of the form 	  $$X=\left[
	 \alpha_1 \quad \alpha_2 \quad \alpha_3J_{1 \times (n_1-1)} \quad \alpha_4 J_{1 \times (n_2-1)}
	 \right]^T$$ for some $(\alpha_1, \alpha_2, \alpha_3, \alpha_4) \neq (0,0,0,0)$.  	 
	 If $\beta$ is an eigenvalue of $AD(S_{n_1,n_2})$ corresponding to the eigenvector $X$, then \[AD(S_{n_1,n_2})X = \beta X, \]	 
%	   $$\left[
%	 \alpha_1 \quad \alpha_2 \quad \alpha_3J_{1 \times (n_1-1)} \quad \alpha_4 J_{1 \times (n_2-1)}
%	 \right]^T,$$ 
%	 
%	 then
%	 \[AD(S_{n_1,n_2}) \left[
%	 \alpha_1 \quad \alpha_2 \quad \alpha_3J_{1 \times (n_1-1)} \quad \alpha_4 J_{1 \times (n_2-1)}
%	 \right]^T = \beta \left[
%	 \alpha_1 \quad \alpha_2 \quad \alpha_3J_{1 \times (n_1-1)} \quad \alpha_4 J_{1 \times (n_2-1)}
%	\right]^T, \]
	 which is equivalent to the system 	 
	 \[\begin{bmatrix}
	 -\beta & 1  & n_1-1 & 0\\
	 1 &  -\beta & 0 & n_2-1\\
	 1 & 0 & -\beta & 3(n_2-1)\\
	 0 & 1 & 3(n_1-1) & -\beta
	 \end{bmatrix} 
	 \begin{bmatrix}
	 \alpha_1 \\ \alpha_2 \\ \alpha_3 \\ \alpha_4
	 \end{bmatrix} = 
	 \begin{bmatrix}
	  0 \\ 0 \\ 0 \\ 0
	 \end{bmatrix}.\] 
	 
	 The above system admits a nontrivial solution if and only if the square matrix on the left-hand side is non-singular. Hence, the remaining four eigenvalues of $AD(S_{n_1,n_2})$ are the eigenvalues of this matrix
	 \[\begin{bmatrix}
	 	0 & 1 & n_1-1 & 0\\
	 	1 & 0 & 0 & n_2-1\\
	 	1 & 0 & 0 & 3(n_2-1)\\
	 	0 & 1 & 3(n_1-1) & 0
	 \end{bmatrix}\]
	 whose characteristic polynomial is $x^4-\left(9n_1n_2-8n_1-8n_2+8\right)x^2+4\left(n_1n_2-n_1-n_2+1\right)$.
	 This concludes the proof.	 
	\end{proof}

\end{document}
